% 02-v3-cfg-typing-classification — manuscript skeleton (started 2026-11-20).
% Source map and design decisions: OUTLINE.md.  Text marked TODO(port: ...) is to be
% copied from the archived 02-v2 sources named there and specialized to fan-out one.
% Abbreviations in source comments:
%   O  = 02-v2-lg-fanout-boundary/parked/planA_2026-11-16/original_en_43pp.tex
%   P1 = 02-v2-lg-fanout-boundary/parked/plan1_small_alphabet_2026-11-11.tex
%   PU = 02-v2-lg-fanout-boundary/parked/parked_u1_typing_level.tex
%   M  = 02-v2-lg-fanout-boundary/main.tex (Plan A, 2026-11-17)
\documentclass[11pt]{article}

\usepackage[T1]{fontenc}
\usepackage{lmodern}
\usepackage{microtype}
\usepackage{amsmath,amssymb,amsthm,mathtools}
\numberwithin{equation}{section}
\usepackage{enumitem}
\usepackage{geometry}
\geometry{margin=1in}
\usepackage{cite}
\usepackage{hyperref}
\usepackage{url}
\usepackage{xcolor}
\hypersetup{hidelinks}
\setlength{\emergencystretch}{3em}

\theoremstyle{plain}
\newtheorem{theorem}{Theorem}[section]
\newtheorem{lemma}[theorem]{Lemma}
\newtheorem{proposition}[theorem]{Proposition}
\newtheorem{corollary}[theorem]{Corollary}
\theoremstyle{definition}
\newtheorem{definition}[theorem]{Definition}
\theoremstyle{remark}
\newtheorem{remark}[theorem]{Remark}
\newtheorem{openproblem}[theorem]{Open Problem}
\theoremstyle{plain}
\newtheorem{mainthm}{Theorem}
\renewcommand{\themainthm}{\Alph{mainthm}}

% Notation (aligned with the companion paper #1 where possible).
\newcommand{\Ccf}[1]{\mathcal{C}^{\mathrm{cf}}_{#1}}
\newcommand{\B}{\mathcal{B}}
\newcommand{\D}{\mathcal{D}}
\newcommand{\eps}{\varepsilon}
\newcommand{\LG}{\mathbf{LG}}
\newcommand{\todo}[1]{\textcolor{red}{\textbf{[TODO: #1]}}}

\title{Which Fixed Typings Give Polynomial Characteristic Data?\\
A Classification for Substitutable Context-Free Languages}
\author{Takayuki Kuriyama\\
\texttt{growup.kuriyama@gmail.com}}
\date{Draft skeleton, \today}

\begin{document}
\maketitle

\begin{abstract}
% DRAFT (2026-11-20). Target <= 200 words.
We study the total length of characteristic samples for positive-data
learning of context-free languages whose substitutability comparisons are
restricted by a fixed morphism $s$ from nonempty words onto a finite
semigroup.  A companion paper shows that such samples are polynomial in
grammar size and thickness when the image of $s$ lies in the join of the
locally trivial semigroups and the finite groups, and that the two-element
semilattice $U_1$ is an obstruction; it claims no necessity.
We determine where polynomial characteristic data ends, at two levels.
At the level of pseudovarieties, for rank cap at least two, polynomial
characteristic data holds uniformly exactly for subpseudovarieties of the
local groups $\LG$; local-group typings admit samples of length
$O(\|G\|^2\bar\theta_G)$, and outside $\LG$ some typing forces exponential
length.  At the level of individual typings onto $U_1$, polynomial data
holds exactly when a single letter has type~$1$; such typings lie outside
$\LG$ and still have polynomial data, but only at a cost quadratic in
thickness.  Finally, the product $\|G\|\bar\theta_G$ is necessary for every
typing, and is attained by the trivial typing.
\end{abstract}

\noindent\textbf{Keywords:} Grammatical inference; Context-free languages;
Distributional learning; Characteristic samples; Finite semigroups; Local groups

%======================================================================
\section{Introduction}
\label{sec:introduction}
% DRAFT. Sources: M Section 1 (setting, Yoshinaka comparison), O Section 1.

\paragraph{Background.}
Distributional learning identifies constituents from positive examples by
their contexts \cite{Gold1967,Angluin1980,ClarkEyraud2007}.  Substitutable
context-free languages are identifiable in the limit from positive data
\cite{ClarkEyraud2007}, and fixed-window variants admit polynomial
characteristic samples \cite{Yoshinaka2008}.  The companion paper
\cite{kuriyamaCFG} replaces windows by a fixed morphism
$s:\Sigma^+\twoheadrightarrow\mathsf S$ onto a finite semigroup: two nonempty
factors are compared only when their $s$-types agree.  It proves exact
reconstruction and Gold identification for every fixed $s$, polynomial
characteristic data in terms of a typed thickness parameter, and polynomial
data in ordinary thickness when $\mathsf S\in\mathbf{LI}\vee\mathbf{Grp}$.
For the two-element semilattice $U_1$ it exhibits targets whose
characteristic samples are exponentially long, and it explicitly leaves open
whether the algebraic hypothesis is necessary.

\paragraph{Question.}
\emph{For which fixed typings is the total length of characteristic samples
bounded polynomially in grammar size and thickness?}  We answer this question
at two levels, and we determine the shape of the bound.

\paragraph{Setting.}
Fix $s$ and a rank cap $r_{\max}$.  The target class $\Ccf{s}$ consists of the
context-free languages in which factors with the same type and a common
accepting context have the same distribution; it coincides with the class of
\cite{kuriyamaCFG} (Lemma~\ref{lem:classes-agree}).  Data length is
$\|K\|_+=\sum_{w\in K}(|w|+1)$, measured against the size $\|G\|$ and the rule
thickness $\bar\theta_G$ of reduced good presentations $G$ of the target.
\todo{one paragraph on the reconstruction operators $\B_s$, $\B_s^{\le r}$ and on
set-driven operators; port from M Section 1 ``Setting''.}

\paragraph{Main results.}
A finite semigroup is a \emph{local group} if $e\mathsf Se$ is a group for every
idempotent $e$; write $\LG$ for this pseudovariety.  It contains
$\mathbf{LI}\vee\mathbf{Grp}$ properly \cite{CostaNogueiraTeixeira2021}.

\begin{mainthm}[Pseudovariety level]
\label{thm:A}
Let $r_{\max}\ge2$ and let $\mathbf V$ be a pseudovariety of finite semigroups.
Then $\mathbf V$ has uniform polynomial characteristic-data length
\textup(Definition~\ref{def:uniform-polydata}\textup) if and only if
$\mathbf V\subseteq\LG$.  More precisely:
\begin{enumerate}[label=\textup{(\roman*)},leftmargin=*]
\item if $s(\Sigma^+)\in\LG$, every target has a characteristic sample of total
length $O(\|G\|^2\bar\theta_G)$, with a constant depending only on $r_{\max}$
and $|s(\Sigma^+)|$ \textup(Theorem~\ref{thm:lg-sufficiency}\textup);
\item if $\mathbf V\not\subseteq\LG$, then $U_1\in\mathbf V$, and a typing onto
$U_1$ of a three-letter alphabet has targets with reduced good rank-two
presentations of size and thickness $O(n)$ whose characteristic samples all
have length at least $3\cdot2^n-1$ \textup(Theorem~\ref{thm:pseudovariety-boundary}\textup).
\end{enumerate}
\end{mainthm}

\begin{mainthm}[Typing level]
\label{thm:B}
Let $r_{\max}\ge2$ and let $s:\Sigma^+\twoheadrightarrow U_1$ be a typing.
Then $s$ admits polynomial characteristic-data length if and only if exactly
one letter has type~$1$ \textup(Corollary~\ref{cor:u1-classification}\textup).
More generally, an isolated letter \textup(Definition~\ref{def:isolated-letter}\textup)
has polynomially short representatives over any alphabet
\textup(Theorem~\ref{thm:isolated-letter}\textup), whereas letter-disjoint words
realizing a copy of $U_1$ force exponential data
\textup(Corollary~\ref{cor:typing-level-obstruction}\textup).
\end{mainthm}

\begin{mainthm}[Shape of the bound]
\label{thm:C}
Let $r_{\max}\ge2$.
\begin{enumerate}[label=\textup{(\roman*)},leftmargin=*]
\item For every typing over an alphabet with at least two letters there are
finite targets, with presentations within a constant factor of the smallest,
whose characteristic samples all have length $\Omega(\|G\|\bar\theta_G)$; for the
trivial typing this bound is attained \textup(Proposition~\ref{prop:linear-factor},
Corollary~\ref{cor:trivial-exact}\textup).
\item For the binary typing $s_{U,2}$ \textup($\mathtt a\mapsto1$, $\mathtt b\mapsto\zeta$\textup),
whose image $U_1$ is not a local group, data are polynomial but not bounded by
$p(\|G\|)\,\bar\theta_G$ for any polynomial $p$: there are targets needing length
greater than $\bar\theta_G^{\,2}/4$ with $\|G\|=O(\log\bar\theta_G)$
\textup(Proposition~\ref{prop:binary-quadratic}\textup).
\end{enumerate}
\end{mainthm}

Thus the two levels genuinely differ: $\LG$ is the exact boundary for
pseudovarieties, while individual typings outside $\LG$ can still be
polynomial, at the price of a superlinear dependence on thickness.

\paragraph{Relation to the companion paper.}
\todo{state precisely what is taken from \cite{kuriyamaCFG} (framework, class,
qualitative reconstruction) and what is new here (LG sufficiency, the two
classifications, the shape results); explain that the companion's
bracketed family and our three-letter family are variants; port the
presentation-relative comparison from M Section 1.}

\paragraph{Proof methods.}
\todo{sliding-factor coordinates for $\LG$ and simultaneous shortening
(Section~\ref{sec:lg}); rare-type and distinguishing-pair lower bounds
(Section~\ref{sec:data}); heavy-path periodicity and counting for isolated
letters (Appendix~\ref{app:isolated}).}

\paragraph{Organization.}
\todo{one paragraph.}

%======================================================================
\section{Preliminaries}
\label{sec:prelim}
% TODO(port: O Section 2 — def:characteristic-sample, def:semigroup-typing,
%   def:h-refinement, def:prefix-suffix-typing, def:standard-mcfg [rewrite for CFG],
%   def:grammar-properties, def:class-cmcf [rewrite as \Ccf{s}]; drop tuple contexts).
\todo{positive-data learning and characteristic samples; typings and the
fresh-identity extension; CFG size $\|G\|$, good and reduced presentations, rule
thickness $\bar\theta_G$; the class $\Ccf{s}$.}

\begin{lemma}[Agreement with the companion class]
\label{lem:classes-agree}
For every typing $s$, the class $\Ccf{s}$ and the notion of characteristic
sample coincide, up to notation, with those of
\cite[Definition~2.1 and Section~3]{kuriyamaCFG}.
\end{lemma}
\begin{proof}
\todo{port: M lem:fanout-one-classes.}
\end{proof}

%======================================================================
\section{Reconstruction and Characteristic Data}
\label{sec:data}
% TODO(port: O Sections 3--4 — def:reconstruction, prop:reconstruction-soundness,
%   prop:completeness, prop:good-normalization, thm:exact-reconstruction,
%   def:witness-set, thm:data-from-typed-thickness, lem:distinguishing,
%   prop:rare-type, prop:typed-thickness-exponential).
% Qualitative parts (soundness, completeness, Gold) are Yoshinaka / companion
% material: cite, do not re-prove at length.
\todo{reconstruction operators, typed witnesses, typed thickness and data length.}

\begin{lemma}[Distinguishing sub-targets]
\label{lem:distinguishing}
\todo{port statement and proof: O lem:distinguishing (f=1).}
\end{lemma}

\begin{proposition}[Rare types force positive data]
\label{prop:rare-type}
\todo{port statement and proof: O prop:rare-type (f=1).}
\end{proposition}

%======================================================================
\section{Local Groups Suffice}
\label{sec:lg}
% TODO(port: O Section 5.1 and 5.3; O app:local-group (thm:lg-characterization,
%   cor:lg-type-level-shortest); O app:group-blocks restricted to one context).
% Follow-up task: replace the multi-block argument by a direct two-block
% (left/right context) argument at fan-out one.

\begin{theorem}[Polynomial characteristic data for local-group typings]
\label{thm:lg-sufficiency}
Let $r_{\max}\ge1$ and suppose $s(\Sigma^+)\in\LG$.  For every $L\in\Ccf{s}$ and
every reduced good presentation $G$ of $L\cap\Sigma^+$ of rank $\le r_{\max}$, the
operators $\B^{\le r_{\max}}_s$ and $\B_s$ have a common characteristic sample
for $L$ of total length $O(\|G\|^2\bar\theta_G)$, with a constant depending only
on $r_{\max}$ and $|s(\Sigma^+)|$.
\end{theorem}
\begin{proof}
\todo{port: O thm:local-group-polydata with $f=1$.}
\end{proof}

%======================================================================
\section{The Pseudovariety Boundary}
\label{sec:pseudovariety}
% TODO(port: O Section 7 — prop:three-letter-family, def:uniform-polydata,
%   thm:fanout-one-data; M rem:fanout-one-data for the scope remark).

\begin{definition}[Uniform polynomial characteristic-data length]
\label{def:uniform-polydata}
Let $r_{\max}\ge1$.  A pseudovariety $\mathbf V$ of finite semigroups has
\emph{uniform polynomial characteristic-data length at rank $r_{\max}$} if for
every nonempty finite alphabet $\Sigma$ and every surjective typing
$s:\Sigma^+\twoheadrightarrow\mathsf S$ with $\mathsf S\in\mathbf V$ there is a
polynomial $p=p_{r_{\max},s,\Sigma}$ such that for every $L\in\Ccf{s}$ and every
reduced good rank-$\le r_{\max}$ presentation $G$ of $L\cap\Sigma^+$, the
operator $\B^{\le r_{\max}}_s$ has a characteristic sample $F$ for $L$ with
$\|F\|_+\le p(\|G\|,\bar\theta_G)$.  A single typing $s$ \emph{admits polynomial
characteristic-data length} if the same holds for that $s$.
\end{definition}

\begin{proposition}[Three-letter doubling family]
\label{prop:three-letter-family}
\todo{port: O prop:three-letter-family; proof in Appendix~\ref{app:three-letter}.}
\end{proposition}

\begin{theorem}[Pseudovariety boundary]
\label{thm:pseudovariety-boundary}
\todo{port: O thm:fanout-one-data, with the operator quantification of M
thm:fanout-one-data(iii).}
\end{theorem}

%======================================================================
\section{Individual Typings}
\label{sec:typings}
% TODO(port: P1 Part 1 — cor:typing-level-obstruction, prop:binary-u1-quadratic,
%   def:isolated-letter, thm:two-letter-fanout-one, cor:two-letter-fanout-one,
%   cor:u1-classification; PU prop:binary-u1-weight-slice for the sharp
%   constant-weight bound).

\begin{corollary}[A typing-level obstruction]
\label{cor:typing-level-obstruction}
\todo{port: P1 cor:typing-level-obstruction.}
\end{corollary}

\begin{definition}[Isolated letter]
\label{def:isolated-letter}
A letter $\mathtt a\in\Sigma$ is \emph{isolated} for a typing $s$ if
$s^{-1}(s(\mathtt a))=\mathtt a^+$ and all words that contain $\mathtt a$ but do
not lie in $\mathtt a^+$ have one common type.
\end{definition}

\begin{theorem}[Isolated letters have short pure representatives]
\label{thm:isolated-letter}
\todo{port: P1 thm:two-letter-fanout-one (bound $C(\|G\|+2)^{22}\bar\theta_G^{\,8}$,
$C<10^{43}$); proof in Appendix~\ref{app:isolated}.}
\end{theorem}

\begin{corollary}[Polynomial data from isolated letters]
\label{cor:isolated-data}
\todo{port: P1 cor:two-letter-fanout-one.}
\end{corollary}

\begin{corollary}[Typings onto $U_1$]
\label{cor:u1-classification}
\todo{port: P1 cor:u1-classification.}
\end{corollary}

\begin{proposition}[Quadratic data for the binary typing]
\label{prop:binary-quadratic}
\todo{port: P1 prop:binary-u1-quadratic and PU prop:binary-u1-weight-slice;
proof in Appendix~\ref{app:u1}.}
\end{proposition}

%======================================================================
\section{The Linear Factor and Rank One}
\label{sec:shape}
% TODO(port: O prop:linear-factor-all-fanout(iv) at f=1 over two letters,
%   O cor:trivial-exact; optional PU thm:rankone-polydata at f=1).

\begin{proposition}[The linear factor is necessary]
\label{prop:linear-factor}
\todo{port: O prop:linear-factor-all-fanout, two-letter fan-out-one case.}
\end{proposition}

\begin{corollary}[The trivial typing]
\label{cor:trivial-exact}
\todo{port: O cor:trivial-exact at $f=1$.}
\end{corollary}

%======================================================================
\section{Conclusion}
\label{sec:conclusion}
\todo{summary; open: classification of individual typings beyond $U_1$
(two non-commuting words of an idempotent type without letter-disjointness);
the factor $\|G\|$ between $\|G\|\bar\theta_G$ and $\|G\|^2\bar\theta_G$;
remark that local-group sufficiency extends to multiple context-free grammars
of every fixed fan-out (not claimed here).}

\section*{Declaration of competing interest}
The author declares no competing interests.

\section*{Data availability}
No data were used.

\appendix

\section{Local-Group Coordinates and Shortening}
\label{app:lg}
\todo{port: O app:local-group and app:group-blocks, restricted to fan-out one.}

\section{The Three-Letter Doubling Family}
\label{app:three-letter}
\todo{port: O app:three-letter-family.}

\section{Isolated Letters}
\label{app:isolated}
\todo{port: P1 app:two-letter-fanout-one, with the 2026-11-09 review corrections
(lem:app-h-tree(d) only for $k\ge1$; empty products and zero exponents explicit).}

\section{Typings onto \texorpdfstring{$U_1$}{U1}}
\label{app:u1}
\todo{port: P1 app:u1-typings.}

\section{The Linear-Factor Family}
\label{app:linear}
\todo{port: O app:linear-factor-fanout at fan-out one.}

\section{Auxiliary Proofs for Typed Witnesses and Normalization}
\label{app:aux}
\todo{port: O app:typed-witness-aux-proofs, app:normalization-details.}

% Bibliography copied from 02-v2-lg-fanout-boundary/main.tex (2026-11-17, Plan A). Prune unused entries before submission.
\begin{thebibliography}{99}

\bibitem{Gold1967}
E.~M.~Gold.
\newblock Language identification in the limit.
\newblock {\em Information and Control}, 10(5):447--474, 1967.

\bibitem{Angluin1980}
D.~Angluin.
\newblock Inductive inference of formal languages from positive data.
\newblock {\em Information and Control}, 45(2):117--135, 1980.

\bibitem{ClarkEyraud2007}
A.~Clark and R.~Eyraud.
\newblock Polynomial identification in the limit of substitutable context-free
languages.
\newblock {\em Journal of Machine Learning Research}, 8:1725--1745, 2007.

\bibitem{deLaHiguera1997}
C.~de~la~Higuera.
\newblock Characteristic sets for polynomial grammatical inference.
\newblock {\em Machine Learning}, 27:125--138, 1997.

\bibitem{Clark2013}
A.~Clark.
\newblock Learning trees from strings: A strong learning algorithm for some context-free grammars.
\newblock \emph{Journal of Machine Learning Research}, 14:3537--3559, 2013.

\bibitem{EyraudHeinzYoshinaka2016}
R.~Eyraud, J.~Heinz, and R.~Yoshinaka.
\newblock Efficiency in the identification in the limit learning paradigm.
\newblock In J.~Heinz and J.~M.~Sempere (eds.), {\em Topics in Grammatical
Inference}, pp.~25--46. Springer, 2016.
\newblock doi:10.1007/978-3-662-48395-4\_2.

\bibitem{Yoshinaka2008}
R.~Yoshinaka.
\newblock Identification in the limit of $k,l$-substitutable context-free
languages.
\newblock In {\em Grammatical Inference: Algorithms and Applications},
LNCS 5278, pp.~266--279. Springer, 2008.

\bibitem{Yoshinaka2009}
R.~Yoshinaka.
\newblock Learning mildly context-sensitive languages with multidimensional
 substitutability from positive data.
\newblock In {\em Algorithmic Learning Theory}, LNCS 5809, pp.~278--292.
 Springer, 2009.
\newblock doi:10.1007/978-3-642-04414-4\_24.

\bibitem{Yoshinaka2011}
R.~Yoshinaka.
\newblock Efficient learning of multiple context-free languages with
multidimensional substitutability from positive data.
\newblock {\em Theoretical Computer Science}, 412(19):1821--1831, 2011.
\newblock doi:10.1016/j.tcs.2010.12.058.

\bibitem{Yoshinaka2015General}
R.~Yoshinaka.
\newblock General perspective on distributionally learnable classes.
\newblock In \emph{Proceedings of the 14th Meeting on the Mathematics
of Language (MoL 2015)}, pp.~87--98.
Association for Computational Linguistics, 2015.
\newblock doi:10.3115/v1/W15-2308.

\bibitem{ClarkYoshinaka2016}
A.~Clark and R.~Yoshinaka.
\newblock Distributional learning of context-free and multiple context-free grammars.
\newblock In J.~Heinz and J.~M.~Sempere (eds.), \emph{Topics in Grammatical
Inference}, pp.~143--172. Springer, 2016.
\newblock doi:10.1007/978-3-662-48395-4\_6.

\bibitem{RambowSatta1999}
O.~Rambow and G.~Satta.
\newblock Independent parallelism in finite copying parallel rewriting systems.
\newblock {\em Theoretical Computer Science}, 223(1--2):87--120, 1999.
\newblock doi:10.1016/S0304-3975(97)00190-4.

\bibitem{KuriyamaMCFGpreprint}
T.~Kuriyama.
\newblock The value of finite observation in positive-data learning
of multiple context-free languages.
\newblock arXiv:2605.11644v3, 2026;
\url{https://arxiv.org/abs/2605.11644v3}.

\bibitem{CostaNogueiraTeixeira2021}
J.~C.~Costa, C.~Nogueira, and M.~L.~Teixeira.
\newblock The word problem for $\kappa$-terms over the
pseudovariety of local groups.
\newblock \emph{Semigroup Forum}, 103:439--468, 2021.
\newblock doi:10.1007/s00233-021-10207-9.

\bibitem{kuriyamaCFG}
T.~Kuriyama.
\newblock Distributional learning of context-free languages under fixed
finite-semigroup typing.
\newblock Manuscript under review for \emph{Theoretical Computer Science};
\url{https://arxiv.org/abs/1409.6247v5}.

\bibitem{Pin1986}
J.-E. Pin.
\newblock {\em Varieties of Formal Languages}.
\newblock North Oxford Academic, London, and Plenum, New York, 1986.

\bibitem{SekiEtAl1991}
H.~Seki, T.~Matsumura, M.~Fujii, and T.~Kasami.
\newblock On multiple context-free grammars.
\newblock {\em Theoretical Computer Science}, 88(2):191--229, 1991.
\newblock doi:10.1016/0304-3975(91)90374-B.

\bibitem{Kallmeyer2010}
L.~Kallmeyer.
\newblock {\em Parsing Beyond Context-Free Grammars}.
\newblock Springer, Berlin, Heidelberg, 2010.
\newblock doi:10.1007/978-3-642-14846-0.

\bibitem{Kanazawa2019Ogden}
M.~Kanazawa.
\newblock Ogden's lemma, multiple context-free grammars, and the control language hierarchy.
\newblock {\em Information and Computation}, 269:104449, 2019.
\newblock doi:10.1016/j.ic.2019.104449.

\bibitem{GomezRodriguezEtAl2009}
C.~G\'omez-Rodr\'iguez, M.~Kuhlmann, G.~Satta, and D.~Weir.
\newblock Optimal reduction of rule length in linear context-free rewriting systems.
\newblock In {\em Proceedings of Human Language Technologies: The 2009 Annual
Conference of the North American Chapter of the Association for Computational
Linguistics}, pp.~539--547. Association for Computational Linguistics, 2009.

\bibitem{OlsonWhite1977}
J.~E.~Olson and E.~T.~White.
\newblock Sums from a sequence of group elements.
\newblock In \emph{Number Theory and Algebra}, pp.~215--222.
Academic Press, New York--London, 1977.

\bibitem{GeroldingerHalterKoch2006}
A.~Geroldinger and F.~Halter-Koch.
\newblock {\em Non-Unique Factorizations: Algebraic, Combinatorial and
Analytic Theory}, volume 278 of {\em Pure and Applied Mathematics}.
\newblock Chapman \& Hall/CRC, Boca Raton, 2006.

\bibitem{AlmeidaEtAl2009}
J.~Almeida, S.~Margolis, B.~Steinberg, and M.~Volkov.
\newblock Representation theory of finite semigroups, semigroup radicals and
formal language theory.
\newblock {\em Transactions of the American Mathematical Society},
361(3):1429--1461, 2009.
\newblock doi:10.1090/S0002-9947-08-04712-0.

\bibitem{Straubing1985}
H.~Straubing.
\newblock Finite semigroup varieties of the form $\mathbf V*\mathbf D$.
\newblock {\em Journal of Pure and Applied Algebra}, 36:53--94, 1985.
\newblock doi:10.1016/0022-4049(85)90062-3.

\bibitem{Tilson1987}
B.~Tilson.
\newblock Categories as algebra: an essential ingredient in the theory of monoids.
\newblock {\em Journal of Pure and Applied Algebra}, 48:83--198, 1987.
\newblock doi:10.1016/0022-4049(87)90108-3.

\end{thebibliography}


\end{document}
